% ============================================================
%  Poul Martin Møller,
%  "Lecture-Paragraphs on Moral Philosophy"
%  English translation of "Forelæsnings-Paragrapher over
%  Moralphilosophien" (Efterladte Skrifter, bd. 5, pp. 141--162).
%
%  Mirrors transcription.tex 1:1. \opage{N} marks the point at
%  which printed page N begins. Numbered lecture-paragraphs (§§).
%  The editor's remarks ("Anm.") are set smaller.
% ============================================================

\documentclass[12pt]{article}
\usepackage[utf8]{inputenc}
\usepackage[T1]{fontenc}
\usepackage{libertinus}
\usepackage{libertinust1math}
\usepackage{textalpha}
\usepackage[english]{babel}
\usepackage[protrusion=true,expansion=true]{microtype}
\usepackage[margin=1.4in]{geometry}
\usepackage{setspace}
\usepackage{fancyhdr}
\usepackage{marginnote}
\usepackage[colorlinks=true,linkcolor=black,urlcolor=black,
  pdftitle={Lecture-Paragraphs on Moral Philosophy},
  pdfauthor={Poul Martin Møller}]{hyperref}
\setlength{\headheight}{15pt}
\pagestyle{fancy}
\fancyhf{}
\fancyhead[LE,RO]{\thepage}
\fancyhead[RE]{\textit{Moral Philosophy}}
\fancyhead[LO]{\textit{Poul Martin Møller}}
\setlength{\parindent}{1.5em}
\setlength{\parskip}{0pt}
\onehalfspacing

\newcommand{\opage}[1]{\marginnote{\footnotesize\textit{[#1]}}}
\newcommand{\ombreak}{\par\medskip\begin{center}\rule{0.32\textwidth}{0.4pt}\end{center}\medskip}
\newcommand{\slaw}[1]{\par\medskip\begin{center}\textbf{\S~#1.}\end{center}\nopagebreak}
\newcommand{\anm}[1]{\par\smallskip{\small\hangindent=1.6em\hangafter=0\textit{Note.}\quad #1\par}\smallskip}
\newcommand{\mchap}[2]{\par\bigskip\begin{center}\large\textbf{#1}\\[0.3em]\textbf{#2}\end{center}\medskip}

\title{\textbf{Lecture-Paragraphs\\[0.2em] on Moral Philosophy}\footnote{From lecture-notes of 1837, collated with the Author's manuscript. F.\ C.\ O.}}
\author{Poul Martin Møller}
\date{Posthumous Writings, vol.\ 5}

\begin{document}
\maketitle

\slaw{1}
\opage{141}According to the common way of conceiving it, morality is the science of man's rational actions and of the aim that is to be attained by them. The cognition of the good which is scientifically developed in moral philosophy is at first present in an immediate way, so that the maxim to be followed in particular cases is discerned without any consciousness of a ground for it. Where the question of the ground of the practical rules enters, it is a natural line of thought to seek the ultimate aim of men's endeavours and from it to derive all the particular rules for them. According to the dogmatic method that thereby arises, the supreme practical rule is expressed in a single formula, which is regarded as self-evident and laid at the foundation of the whole system. In this way several opposed principles have, in the course of the science's historical development, been asserted against one another, and the most important duties and virtues have with greater or lesser one-sidedness been developed out of each of them, so that none of them can be wholly wrong. But by subjecting them to a scientific critique one brought their one-sidedness to consciousness and reached the result that man's rational striving cannot be determined by any single universally valid rule. In the properly philosophical method these two views are united, in that the defect in every one-sided principle is demonstrated, but thereby at the same time a transition is made to a new determination of the principle, so that the true principle of the science is de\opage{142}veloped by the thinking that passes through all the one-sided determinations and follows their natural transition into one another.

\slaw{2}
In morality it is presupposed that man is a free being, but the concept of freedom receives new determinations within the science itself. --- The freedom that is presupposed is not that assumed by the indifferentists, who maintain that man can, in every moment, determine himself against all motives. These defend their theory partly by an alleged immediate consciousness that one does sometimes really determine oneself against all motives; partly by the assertion that man, when the motives are equally strong, cannot determine himself except with absolute arbitrariness; they hold, finally, that the particular action can, only on their view, have moral worth or unworth. But to this it must be answered, as to the first, that immediate experience cannot teach anything negative; as to the second, that the presupposed absolute equilibrium cannot occur; as to the third, that the action can be judged as separate from the character only by an incorrect abstraction. Moreover indifferentism, pursued in its consequences, is wholly unthinkable and, in abolishing the concept of character, conflicts with the express testimony of experience.

\slaw{3}
On the other side, one-sided determinism is also incorrect, when it either makes the will wholly dependent on the scientific development of cognition, or when it assumes that the particular action is a necessary product of the character and the outward circumstances, and the character a result of all the preceding actions. If it then regards the first action as free (predeterminism), it is herein inconsistent; for the essential difficulties that occur in comprehending the freedom of the subsequent actions apply just as fully to the first action. If, on the contrary, the first action too is taken to be unfree, \opage{143}then determinism is not essentially different from fatalism, which conflicts both with moral consciousness and with the genuine concept of a creating Godhead.

\slaw{4}
According to the correct concept of freedom, the particular action is thought neither to be quite independent of \emph{the empirical character}, which would abolish the concept of virtue and vice, nor to be wholly determined by it; on the contrary, the empirical character can in a moral respect be altered, yet in such a way that the altered direction of the will only gradually brings about the facility [Færdighed] in the good or in the evil. An immoral character can never be wholly unalterable, but a virtuous character can, on the other hand, be so consolidated that no alteration of principles can take place in it. --- Motives determine the actions, but they are never merely outward. Only regarded from the one side can they be merely outward; from the other they are inward and come to hold good through the will's determinate direction. Conversion may indeed occur suddenly, in so far as altered maxims can be followed at one stroke; but the total alteration of the disposition comes about only through much exercise.

\slaw{5}
Moral significance belongs only to those expressions of life which man has been aware of and which he has willed. The various degrees of freedom that may here be thought to occur are developed in the first main part of morality, \emph{the Doctrine of Imputation}, in which the concept of an action is determined without regard yet being taken to what it aims at; in the second main part, \emph{the Doctrine of Happiness}, the various purposes of human striving are developed, and, since these come to conflict with one another, but must reasonably be subordinated to a universal aim, there is required yet a third main part of morality, \emph{the Doctrine of the True Good}.

\mchap{First Chapter.}{On Imputation.}

\slaw{6}
\opage{144}Man is said to perform an action when he is the voluntary cause of a change in the surrounding world. In the inviolable determination of an action, voluntariness, two determinations are to be distinguished: 1) \emph{self-determination under the exerted activity}, 2) \emph{the apprehension of the objects} at which it aims; and in the merely voluntary action these two determinations coincide in a single moment. If one of these determinations is lacking, so that man is moved by an irresistible outward power, or has incorrect representations of the objects on which he acts, then the effects of which man is the cause are to be regarded as involuntary. A combination of the voluntary and the involuntary occurs in the mixed deeds, which are half voluntary and half involuntary. These occur \emph{partly} when man is forced to perform an action which he does not will in and for itself, but only in order to avoid what in his opinion is a greater evil, \emph{partly} when he finds himself in a natural condition that brings with it a darkened consciousness. --- In so far as all objects of the action have in themselves a constitution that contributes to determining the constitution of the action, properly every human action is mixed.

\anm{Among the errors that make men's actions involuntary is not counted ignorance of the opposition between good and evil in general.}

\slaw{7}
\opage{145}When man makes the experience that the outward circumstances may be apprehended incorrectly and may hinder the execution of the will, his consciousness is developed in such a way that he first forms a resolve [Forsæt], in which the action has an inner existence in his consciousness, and only thereafter carries it out in the world of sense. In the intentional action three sides may be distinguished: \emph{deliberation, decision, and execution}. Every intentional action is voluntary, but an action can be merely voluntary without being intentional: when it is intentional, it is more imputable than as merely voluntary.

\slaw{8}
But the harmony we have supposed to occur between the resolve and the execution of the action can be disturbed by the fact that the apprehended circumstances may be connected with others not apprehended, which either alter the action so that it does not answer to the resolve, or give it unforeseen consequences. When the alteration of the action or the unforeseen consequences could not or ought not to have been foreseen, there occurs a \emph{chance} [Tilfælde] which cannot be imputed at all.

\slaw{9}
When man acts without perceiving the possibility that by his action he may become guilty of an event, although he could and ought to have foreseen it, then the unintentional is imputed to him, yet only as \emph{negligence} [Forseelse]. This has its ground in a neglect of self-activity, which draws after it an ignorance or natural condition, by which the particular error in consciousness that is always connected with the negligence is caused. Most immediately, then, the ignorance or natural condition is imputable, and mediately the action occasioned thereby.

\slaw{10}
More imputable than the merely intentional action is the \emph{deliberate} [overveiede], to which one may in general reckon the action carried out according to plan. This clearly occurs when an action is not merely prepared by several preceding ones, by which the agent has put himself in a position to carry it out, but the series of these actions has been interrupted by intervening actions with other intentions.

\slaw{11}
\opage{146}To action in the widest sense of the word one also reckons the spiritual expressions of life, representations and thoughts. These may sometimes arise involuntarily in man according to the laws of association of ideas, but even the most involuntary actions are nonetheless dependent on man's character. Still more does it depend on man's freedom whether he will interrupt or continue the train of representations that accidentally presses upon him. The voluntary reproduction of representations, fantasizing and thinking, can also, like action in the narrower sense or outward action, be undertaken both according to resolve and according to plan.

\slaw{12}
In every intentional action man always has a \emph{final intention} or a \emph{purpose in view}: he never performs an action merely in order to perform it. The nearest intention in an action may indeed be another action, which he by the first will make possible; in the second action he may have a third as his intention, and so on; but the series of actions always aims, in the end, at a purpose. These purposes constitute the merely inner side of the actions, and with their investigation we enter upon the proper domain of morality.

\anm{The consciousness that man's actions are always means has given occasion to the erroneous view that man's will should be judged merely according to the purpose of his actions. No action can be regarded \emph{merely} as a means: the sum of all a man's actions is the objective image of his will.}

\mchap{Second Chapter.}{On Happiness.}

\slaw{13}
Of the content of the right purpose we know, for the first, nothing but that there must be a purpose for man's en\opage{147}deavours, i.e.\ that man must seek personal satisfaction in his actions or have \emph{interest} in them. Herein it is still left wholly undetermined of what kind the interest is (whether it aims, e.g., at the satisfaction of sensuous desire or at the attainment of a rational aim). This concept of interest, regarded as the supreme moral principle, is the most contentless principle of all, yet it does bring out a necessary determination in the good.\footnote{The Author's manuscript: at man's aim. F.\ C.\ O.}

\anm{In so far as satisfaction, regarded as the aim of the actions, presupposes that one want after another must be filled, \emph{activity}, as the abiding element in these transitions, can also, by a one-sided abstraction, be regarded as the adequate expression of man's vocation.}

\slaw{14}
By the fact that man in his actions carries out his interest or the aim of his life, a harmony is brought about between his will and the world of objects, and the most immediate way in which this agreement comes to the agent's consciousness is through \emph{feeling}. The \emph{agreeable feeling} one is therefore, by a natural line of thought, brought to regard as the purpose of human activity, and this now becomes the object of our consideration. The subject's harmony with the surrounding world is thought as something that is to be brought about and therefore does not yet occur. Man, in his relation to the surrounding world, does not correspond to his vocation: he has it in view, as an objective purpose; but his agreement or disagreement with it we as yet know only as the agreeable or the disagreeable. The principle thus becomes: to shun the disagreeable, or, what is the same, to strive after the agreeable.

\slaw{15}
Among the moralists who place the genuine pleasure of life in the satis\opage{148}faction of the natural desires, the \emph{Cyrenaic school} comes forward with the principle of concentrating the greatest possible pleasure in each single moment of life and of placing man's vocation in an unbroken series of particular enjoyments. But this rule abolishes itself at the slightest reflection, since this leads to taking the consequences of pleasure into consideration and sometimes to denying oneself a pleasure in order to avoid a pain or to attain a greater pleasure, i.e.\ to transgress the principle for the principle's sake. But once one has granted \emph{thought} its right and perceived that it leads man, out of prudence, sometimes to submit to a displeasure, then it is easily recognised that it does not stand in the power of mortals to procure for themselves a preponderant sum of agreeable moments, and it is consistent to pass over thence to the \emph{Cynic principle}: to secure happiness by an independence attained through the satisfaction of as few desires as possible. These two opposed directions were united by Epicurus, in that he indeed declared pleasure to be the highest purpose, yet partly saw the necessity of self-denial, partly began to recognise the significance of spiritual pleasure. \emph{Spiritual pleasure}, which is the true human pleasure, cannot be placed in the momentary satisfaction of the mutually conflicting drives, but only in the execution of a single universal purpose, whereby unity is brought into human activity. Since the activity that promotes the right aim must be inseparably connected with a pleasure, the \emph{eudaimonistic principle} too, rightly understood, contains one of the most important determinations of the moral principle. But to find the general direction of the will that brings pleasure with it, we have as yet no other guidance than \emph{feeling}, which at this stage shows itself as \emph{moral feeling}.

\slaw{16}
When \emph{moral feeling} is assumed to be the right principle of the actions, it is not the momentary satisfaction of the \opage{149}sensuous desires that is regarded as the source of pleasure, but a general direction of the will is assumed to be the right one, because it in and for itself pleases the individual, whether he attains his purpose or not. In this principle lies the truth that the genuine moral will must be so much an expression of the innermost essence of personality that man finds himself well in its demands. Herein lies that true morality is inseparably connected with a peculiar pleasure. True morality must consequently also have the form of feeling, but the content of this feeling is not yet known. This last circumstance has led others to regard the content as being given wholly from without, through \emph{upbringing and the institutions and statutes of the State}. But by this opinion, which really only pushes the solution of the problem further out, the moral freedom is wholly annihilated, according to which man himself can determine the aim of his life. --- If it is asked more closely about the abiding and universal element in that direction of the will which pleases man, then it can, for the first, safely be regarded as an endeavour to preserve this own I; and many have therefore assumed \emph{self-preservation} to be the aim of life.

\slaw{17}
\emph{Self-preservation} as the principle of man's actions brings him to seek pleasure as a means for self-preservation. It is regarded at this standpoint as the aim of life to bring about and make use of such a relation between the subject and the surrounding world that everything outward serves the furtherance of the individual life. In its lowest shape self-preservation is \emph{physical self-maintenance} and shows itself in this, that man strives to secure his bodily existence against attack from without and to procure for himself nourishment and stimulus for the continuation of the merely organic life. Next it shows itself as a striving after \emph{possession} of the natural things which are conditions of his physical existence, or a striving after \opage{150}property. Through the striving after property man comes to stand in relation not merely to the world of sense, but to other men, who acknowledge his right to what he possesses. The acknowledgement of the subject's will that occurs in property is the beginning of his outward \emph{honour} or his esteem among other men, which can also for itself be made a purpose. (A consistent continuation of the striving after honour is the striving after \emph{freedom} as well as the striving to \emph{rule}.) In so far as the subject, in his craving for honour, ascribes to other rational beings an independent worth, there lies herein an element which, when developed, also brings man to set before himself \emph{the welfare of others} as the aim of his actions.

\anm{Although one cannot, with sundry, especially French, philosophers, derive all man's duties and virtues from the striving after self-preservation as here explained, self-preservation in a higher sense is nonetheless an essential determination of the rational disposition.}

\slaw{18}
\emph{Benevolence}, according to which man finds pleasure in the welfare of other men and wishes their furtherance, first comes to proper actuality in friendship and the kindred \emph{social relations}, in which the benevolence is mutual. In its least developed shape friendship is grounded merely on the pleasure that the one finds in the other's company. The more prudent friendship, in which the mutual outward advantage is the ground of the friendship, is, where no other ground occurs, only improperly to be called friendship. But in a third kind, the rational and true friendship, in which a complete harmony between several free wills occurs, both the agreeable and the useful are preserved. By the fact that this relation (true \emph{sociability})\footnote{The Author's manuscript: (of sociability). F.\ C.\ O.}, in which self-preservation and the endeavour for others' welfare are brought into connection with one another, \opage{151}is assumed to be man's vocation, a considerable advance is made in moral consciousness; but when this connection is regarded merely as a consequence of natural inclination, it is accidental and has partiality as its result. The true reconciliation between self-preservation and benevolence first occurs in \emph{right} [Retten], in which the person counts not as friend, but as person, and each receives what is due to him.

\slaw{19}
\emph{Right}, the true rule for the distribution of outward goods in human society, is a purpose that is different both from the solitary and from the social: we call it an \emph{ideal purpose}. The consciousness of it is the emerging consciousness of the Idea of the good, but of the Idea under a one-sided shape. Strict right may so conflict with the demands of benevolence that its enforcement becomes wrong, that strict right becomes injustice and must yield to \emph{equity}. --- The cognition that is necessary in order that right may be brought to consciousness may also come forward as a peculiar \emph{ideal purpose}, which may also be taken to be man's true vocation; man then lives for \emph{the cognition and presentation of the true and the beautiful}, for science and art.

\slaw{20}
Cognition can also, for itself, be made one-sidedly into a purpose, as among those who regard the \emph{contemplative life} as man's vocation. The aim here is \emph{partly} scientific cognition, which is developed both in abstracted concepts and in pure concepts of reason, \emph{partly} aesthetic intuition of the Ideas that lie at the foundation of existing things and which, through a higher activity of the imagination, not determined by experience, are presented by the artist or by those whose imagination is set in motion by his works. When the one-sided theoretical direction is brought into connection with the above-mentioned practical purposes (the solitary, \opage{152}the social, and the purposes of right), the right aim is found to be \emph{action according to cognition}.

\slaw{21}
To act according to cognition belongs not merely, in every moment, following a maxim by which the aim of life is promoted, but also having general knowledge of the nature of existing things, and the power of judgement to apply the knowledge, according to the maxim, in occurring cases. Moreover a firm character is required, in order not, out of regard to the agreeable or disagreeable, to let oneself be carried away to act against cognition. But the aim, which is the main thing, becomes right only by this, that the inner harmony occurs between man's lower and higher inclinations. From this springs also agreement with the outward order, or the will to conform to the nature of objects and their relation to one another. Action according to cognition thus develops into a \emph{striving after harmony among all man's purposes} and, as a consequence thereof, among all his outward relations. He who strives, in a harmonious way, to realise all these purposes, among which cognition itself must be reckoned, strives after the good or the \emph{perfect existence}.

\slaw{22}
\emph{Perfection}, or perfect existence, is a harmonious development, determined by reason, of all man's capacities and relations. To regard perfection as a purpose which for man is absolutely unattainable, from which it follows that not perfection itself, but the approximation to it continued to infinity, must be assumed to be man's vocation, is an error. In every moment in which reason determines man's manner of action, a passing appearance of perfection occurs, and the concept of the infinite progress has reality only in so far \opage{153}as the series of mortals' perfect expressions of life here in life cannot become perfectly continuous.

\mchap{Third Chapter.}{On the Good.}

\begin{center}\textbf{A.\ On Virtue.}\end{center}

\slaw{23}
The \emph{perfect life}, which in the foregoing was recognised to be man's vocation and is the same as the good, consists in \emph{the practical activity of reason}, for which \emph{the natural purposes} are means; but since the good can come to actuality only through them, they are themselves \emph{relative goods}. To these belong both the \emph{bodily goods} (health, strength, agility) and the \emph{goods of outward fortune} (property, honour, etc.); these constitute, together with the inward perfection --- which cannot occur where these goods are wholly lacking --- the \emph{complete good}. But the \emph{highest good} is the inward perfection, or virtue proper.

\slaw{24}
\emph{Virtue} is the inner condition in which all man's inclinations are brought under the dominion of reason. The more virtue is true virtue, the more has the higher power of reason so penetrated the natural inclination that \emph{virtue has become a second nature to man}, or that the rational actions accord with his inclination. But this perfection of the disposition is a facility acquired through exercise, and not any merely natural condition.

\slaw{25}
The rational cognition which is one side of virtue may indeed be promoted by the instruction of others; but \opage{154}it depends, in the end --- like all cognition --- on man's free self-activity, and this all the more as the good cannot be clearly cognised except by him who has himself freely practised it. --- The exercise in virtue that is required in order that it may become a facility may also be promoted by discipline, in so far as men can be compelled by others to arrange their outward actions according to the law of reason. But this outward action, in agreement with reason, can only be a means of furtherance for the development of virtue. Virtue itself, as inward perfection, can be acquired only through the subject's free self-activity.

\slaw{26}
Virtue is no merely \emph{dormant capacity}, but an \emph{active force} that lays itself open to view in outward actions. Therefore no one can know his own or others' virtue except from the actions in which it comes forth. All other supposed consciousness of one's own virtue is partly \emph{self-deception}, partly \emph{moral fantasizing}.

\slaw{27}
In so far as virtue is not wholly opposed to man's natural condition, but is the determination of that condition by reason, it shows itself, in its outward appearance, as a \emph{mean between two extremes}, in that every natural inclination and its outward expression may, in comparison with the demands of reason, be too strong or too weak. But it is only in its outward appearance, in inclinations and actions, that virtue can be regarded as limited by measure. In its essence it is a single force which cannot be too great.

\slaw{28}
The mean between two extremes spoken of in the preceding paragraph is not the same for all. Here especially the different \emph{individuality} comes into consideration, which makes one en\opage{155}deavour more fitting for one, another for another. Therefore some have also with reason wished to divide the virtues according to the different shapes which virtue, according to the difference of capacities, must assume in different individuals. Thus one may in general distinguish the \emph{contemplative} and the \emph{practical characters}, and since these are distinguished merely by the different main direction, the same division lets itself be carried through on both sides.

\slaw{29}
\emph{Misapprehension of the nature of individuality}, or deliberate deviation from its guidance, produces \emph{untruth} in man's life, whereby this is more or less transformed into a life of mere show. Here \emph{affectation} is to be noted, which occurs when man, in order to attain one of the natural purposes, displays a different relation between his person and things than the one that really obtains. He is thereby brought to express assumed feelings, representations and inclinations which he imagines himself to have. With regard to the greater or lesser connection which affectation has with the character, it can show itself now as \emph{momentary}, now as \emph{fixed} (in that a feigned trait is taken up into the character), now as a \emph{habit of affectation}, which according to circumstances changes its appearance.

\slaw{30}
The virtuous man possesses the proper \emph{freedom}, since the cognition according to which he acts is not merely a knowledge of the circumstances, which can also occur in one who follows his drives, but a \emph{rational cognition} that determines the aim of his actions. Such a cognition it stands in \emph{man's own power} to develop in himself, and it is necessarily developed in him who does not, by the habit of following the agreeable instead of the cognised good, stifle its emerging germs in himself.

\anm{On moral weakness and badness.}

\slaw{31}
\begin{center}\emph{Virtue in its outward appearance.}\end{center}
\opage{156}We consider first virtue in its relation to the solitary drives. Since man by nature shuns the sensuously disagreeable and strives after the sensuously agreeable, virtue shows itself here partly as \emph{courage}, partly as \emph{temperance}. Courage is virtue as readiness to go to meet displeasure and dangers when reason requires it. It lies between \emph{cowardice} on the one side and \emph{foolhardiness} on the other. Temperance is the rational measure in the striving after sensuous pleasure. The vice opposed to it on the one side as irrational excess is \emph{intemperance}; on the other side there is an extreme in the practice of depriving oneself of the sensuous enjoyments of life out of misunderstood asceticism or avarice.

\slaw{32}
Both in relation to the natural shunning of the disagreeable and to the craving for the agreeable, there shows itself partly the rational \emph{perseverance}, which consists in this, that man, without being drawn aside by those two hindrances (pleasure and displeasure), pursues a planned rational activity; partly the rational \emph{energy} [Driftighed], which is laid open to view by the fact that man's actions, in proportion to his natural endowment, follow quickly upon one another. Combination of perseverance and energy gives \emph{industriousness}, or planned energy, a virtue that lies between \emph{indolence} and \emph{self-torment}.

\slaw{33}
In its relation to the natural craving for property, virtue shows itself as \emph{good husbandry} [Huusholderiskhed], which is the rational moderation in acquisitiveness and expenditure, and has its origin in correct concepts of property. It lies between \emph{avarice} (covetousness and stinginess) on the one side, which comes from ascribing to property an absolute value, and \emph{prodigality} on the other \opage{157}side, according to which one makes such great expenditures for less important aims that one lacks what is needful for the attainment of important aims.

\slaw{34}
In the drive for honour, the consciousness of one's own worth, or \emph{self-feeling}, lies at the foundation of man's craving that others should acknowledge it. The rational degree of self-feeling, which can very well coexist with humility, lies between \emph{pride}, under which the representation of one's own worth above others too much predominates, and \emph{pusillanimity}, which brings man irrationally to limit his striving and his demands. The rational disposition with regard to honour is called \emph{love of honour} [Ærekjærhed] and is equally far removed from \emph{ambition} and from a lack of the sense of honour.

\slaw{35}
From love of honour springs \emph{love of freedom}, according to which man strives after a condition in which he, without being compelled by other men, can himself determine his activity; a virtue that lies between \emph{licentiousness} and \emph{servility}. A more positive form is taken by the \emph{desire for influence} or \emph{the love of ruling} that springs from the same source, which is rationally so limited that it does not rise to lust for domination, yet is so vigorous that it does not sink down to indifference toward the common affairs of society.

\slaw{36}
With regard to the feelings that arise at the representation of others' condition, virtue shows itself as the \emph{sympathetic mind}, which lies between \emph{hard-heartedness} and \emph{philanthropic sentimentality}. With regard to the need to find sympathy in others, the rational disposition shows itself as \emph{openness}, which on the one side is opposed to \emph{loose-tonguedness}, on the other side to excessive \emph{reserve}. Under the \opage{158}social connections that have merely pleasure as their purpose, virtue comes forth as \emph{affability}, which lies between \emph{complaisance} and \emph{brusqueness}. But in the connections where usefulness is the purpose, virtue as \emph{obligingness} holds the right mean between \emph{selfishness} and self-sacrifice for others carried too far.

\slaw{37}
In his serious striving, the rational man sees himself not merely compelled to counteract the plans of the bad, but sometimes becomes at variance with the better about the means to be chosen for the furtherance of the common purposes; then the right mean shows itself partly as \emph{manliness} or manly mind, in that one makes one's rational will prevail without being hindered by love of ease or by fear, partly as \emph{compliance}, in that one conforms to others' will where one's own principles permit it. The corresponding excesses are partly \emph{lust for domination} and \emph{irascibility}\footnote{The Author's manuscript has: partly irascibility and obstinacy. F.\ C.\ O.}, partly weak \emph{yieldingness} and the sacrifice of one's individual character.

\slaw{38 and 39}
1) In relation to the drive for right, virtue shows itself as \emph{equity of mind} [Retsindighed]. It holds the right mean between abstract \emph{justice} on the one side, which, without regard to the peculiar circumstances that transform strict right into injustice, defends the dead rule of right, and individual \emph{arbitrariness} on the other side, which strives too much to make the subjective good will prevail against definitely enunciated statutes and laws.

2) The right \emph{love of art and science} lies between the lack of any sense for their Ideas and the \opage{159}idolatry of art and science that would wholly separate them from the practical purposes of life, which are merely their necessary presuppositions.

\slaw{40}
Since it is reason that determines the right measure in inclinations and actions, the good is not good for the drive's sake, but because it is sanctioned by reason. By the consciousness that reason, in opposition to the drive, is the higher and more essential side, man is brought to regard the mere satisfaction of natural desires as \emph{worthless} (Stoicism), and to set store merely by the \emph{independent thinking} by which man feels himself free and independent of the lower purposes. Once this standpoint is attained, the practical purposes cannot be recognised to be goods without their content being determined by pure reason instead of by the drive. This demand has, especially in the Christian era, been made of moral philosophy (the doctrine of duty).

\begin{center}\textbf{B.\ On Duty.}\end{center}

\slaw{41}
At the standpoint of the doctrine of duty it is recognised that reason alone shall determine the content of the purpose for man's actions, and that these ought not to be undertaken in order to satisfy the natural desires (heteronomy), but out of \emph{obedience to reason's higher law}, to which man shall unconditionally submit himself (categorical imperative). This consciousness is opposed to the ancient doctrine of virtue, according to which the merely natural desire, within a certain limit, was good and furnished the material regulated by reason. (Although he who acts out of reverence for duty submits himself to a strict command that has the character of a higher power, the freedom of the will is nonetheless saved by the consciousness that reason is man's own innermost essence.) \opage{160}When this view is consistently pursued, the activity of reason must be recognised as the sole purpose that ought to be craved for its own sake, and the natural purposes must be regarded as \emph{relative goods}, which have worth only as means for the furtherance of this purpose.

\anm{On the division of duties. Unconditional duties, indifferent actions.}

\slaw{42}
\emph{Collision of duties} occurs when reason seems to enjoin several actions that stand in such a relation to one another that the execution of the one makes the other impossible. Even though, for the collision of duties, several \emph{rules of prudence} may be given, by which the number of collisions will be diminished, yet no universally valid rules can be given by which any ranking-order [Rangforordning], valid for all cases, among the general classes of duties can be determined. Of the three main classes, neither the duties toward ourselves nor toward others nor toward God can be set up as those which in all cases are to be preferred. Just as little can a universally valid order be determined for the subordinate duties within any single one of these main classes. The understanding must then give up its attempt to give the doctrine of duty its completion in detail, and there remains only the one expedient, of working at the \emph{overcoming of sensuousness}, which hinders the correct application of the commands of duty in the particular cases.

\slaw{43}
By \emph{ascetic exercises} man seeks to \emph{overcome sensuousness} as a hindrance to the activity of reason. \emph{Penitential exercise} and the \emph{self-examination} conducted in solitude are used as means to acquire virtue; but asceticism especially recommends continually directing thought in prayer to the Godhead --- the source of duties. By these means of virtue, especially by the last and most perfect, man strives to bring about \opage{161}such an agreement between himself and the law of duty that he, of himself, as a single individual, can draw the discernments of the command of duty. He then acts according to the promptings of his own heart, or according to his conscience.

\begin{center}\textbf{C.\ On Conscience.}\end{center}

\slaw{44}
\emph{Conscience} cognises immediately which actions are enjoined by reason, and the subject is, in virtue of it, wholly certain that in his action he is in agreement with the good (he has a good conscience). It even sets itself above the definite prescriptions of the moral law, because in the particular case it is conscious of having the right to decide their relative validity. --- But conscience can nonetheless be wavering, indeed erring, since it can even refute itself, and it must consequently come down to the content of the action whether it is to be judged good or not. Thus conscience can, by making itself into the source of the law, come to determine itself against the moral law, or toward evil.

\slaw{45}
Evil can be practised in such a way that the subject has consciousness of the command of duty but nonetheless acts against it; then he is said to act with an \emph{evil conscience}; but he can also, with \emph{hypocrisy} toward himself and others, gloss over his evil action by letting some ground of motivation or other count as an acceptable and good ground for it (\emph{probabilism}). Since such an acceptable ground can be given for every conceivable action, there remains nothing of specially good content, but only the general good intention.

\slaw{46}
The good intention is then assumed to sanctify the means, and \opage{162}since such an intention can be brought out in every bad action, the supposed good is thereby wholly transformed into evil. A purpose is chosen because it appears good to the subject. He acts according to his opinion, and it makes no difference to him whether that which according to his opinion is good is in itself evil. Since, with this reliance on subjective opinion, the good in and for itself is denied, there is thence an easy transition to \emph{moral irony}, according to which one holds one's individual will to be the highest and consequently believes oneself raised above the moral law. Herewith subjective freedom has attained its formal completeness, but lost all objective content.

\slaw{47}
It has been shown in the last paragraphs that the endeavour to cognise man's vocation merely from the individual's standpoint is fruitless and can, in its one-sidedness, lead to principles that conflict with reason. The good shall not be cognised merely by subjective thinking, but is found in its actuality and development in the State, and the individual fulfils his vocation only by acting as a member [Led] in such a rationally organised society. Thereby the individual freedom is not abolished or restricted; for the individual must be able to find again the same reason that develops itself in his own self, in the State's laws, statutes and institutions. He finds therein only a limit for his arbitrariness, in that the fundamental forms of rational freedom, which the true conscience cannot bid man overstep, have therein their objective existence. In the life of the State the satisfaction of the natural drives receives its ethical significance, in that they become means for the State's life and progress, and the freest development of the individual capacities furthers the perfection and happiness of society.

\ombreak

\end{document}
